\begin{abstract}

Audio classification is a fundamental task pivotal to intelligent real-world systems, including voice control and environmental monitoring. However, conventional Convolutional Neural Networks (CNNs) often suffer from large parameter redundancy and limited generalization capability, whereas modern Transformer-based architectures typically necessitate significant computational overhead.
In this study, we present an empirical investigation of the ConvMixer architecture—a simple yet powerful model originally developed for computer vision—for audio classification using the Mel-spectrogram representation.
To our knowledge, this work represents an early and systematic effort to adapt ConvMixer to the audio domain, bridging vision-inspired patch-based design and efficient speech understanding.
The proposed framework incorporates a robust signal processing pipeline (denoising, WebRTC VAD, normalization) to generate fixed-size (128×32) Mel-spectrogram inputs. The ConvMixer model was configured with 8 deep blocks and 256 channels, utilizing approximately 0.7 million parameters.
On a Vietnamese speech command dataset of 1,853 samples and 13 classes, the ConvMixer-256/8 configuration achieved strong representative single-run performance (test accuracy 97.42\%, macro/weighted F1-score 0.97). To improve fairness and statistical reliability, we further define a controlled benchmark protocol in which ConvMixer-256/8, ResNet-18, MobileNetV2, EfficientNet, and AST (tiny configuration) are trained under the same preprocessing, augmentation, split policy, and training budget, and evaluated across five random seeds (42--46) using mean$\pm$std reporting. Under this controlled five-seed setting, ConvMixer-256/8 reached 97.63$\pm$0.59 test accuracy, with macro-F1 0.9761$\pm$0.0059 and weighted-F1 0.9763$\pm$0.0059, while using only 0.72M parameters.
Overall, these results and the revised protocol indicate that a purely convolutional patch-based architecture is a parameter-efficient and practical alternative for real-time speech command recognition, with clear potential for embedded and edge-AI deployment.

\end{abstract}
